\documentclass[12pt, a4paper,
  %oneside,      %% -- odkomentujte, pokud chcete svou práci mít pouze jednostrannou, mezera pro hřbet pak automaticky bude pouze na levé straně
 twoside,        %% -- pro oboustranné práce, mezera pro hřbet následně střídá strany.
 openright
]{report}

%% Nutné balíčky a nastavení
%%%%%%%%%%%%%%%%%%%%%%%%%%%%

\title{Codeventure} %% -- Název tvé práce
\author{Michal Maděra, Josef Kahoun} %% -- tvé jméno
\date{2022} %% -- rok, kdy píšeš SOČku

\usepackage[top=2.5cm, bottom=2.5cm, left=3.5cm, right=1.5cm]{geometry} %% nastaví okraje, left -- vnitřní okraj, right -- vnější okraj

\usepackage[czech]{babel} %% balík babel pro sazbu v češtině
\usepackage[utf8]{inputenc} %% balíky pro kódování textu
\usepackage[T1]{fontenc}
\usepackage{cmap} %% balíček zajišťující, že vytvořené PDF bude prohledávatelné a kopírovatelné

\usepackage{graphicx} %% balík pro vkládání obrázků

\usepackage{subcaption} %% balíček pro vkládání podobrázků

\usepackage{hyperref} %% balíček, který v PDF vytváří odkazy

\linespread{1.15} %% řádkování

\usepackage[pagestyles, nobottomlines]{titlesec} %% balíček pro úpravu stylu kapitol a sekcí
\titleformat{\chapter}[block]{\scshape\bfseries\LARGE}{\thechapter}{10pt}{\vspace{0pt}}[\vspace{-22pt}]
\titleformat{\section}[block]{\scshape\bfseries\Large}{\thesection}{10pt}{\vspace{0pt}}
\titleformat{\subsection}[block]{\bfseries\large}{\thesubsection}{10pt}{\vspace{0pt}}

\setlength{\parindent}{0em}
\setlength{\parskip}{1em}

\setcounter{secnumdepth}{4}
\setcounter{tocdepth}{1}
\usepackage{fancyhdr}
\pagestyle{fancy}
\renewcommand{\headrulewidth}{1pt}

\usepackage{booktabs}

\usepackage{url}

%% Balíčky co se můžou hodit :) 
%%%%%%%%%%%%%%%%%%%%%%%%%%%%%%%

\usepackage{pdfpages} %% Balíček umožňující vkládat stránky z PDF souborů, 

\newcommand{\subsubsubsection}[1]{\paragraph{#1}\mbox{}\\}


%% Začátek dokumentu
%%%%%%%%%%%%%%%%%%%%
\begin{document}

\pagestyle{empty}
\pagenumbering{Roman}

\begin{titlepage}
    \bfseries{
        \begin{center}
            \LARGE{STŘEDOŠKOLSKÁ ODBORNÁ ČINNOST}

            \vspace{14pt}
            \large{
                Obor č. 18: Informatika
            }

            \vspace{0.35 \textheight}

            \LARGE{
			Codeventure
            }

            \vspace{0.35\textheight}
        \end{center}
        
        \noindent\Large{Michal Maděra, Josef Kahoun}

        \noindent\Large{Pardubický kraj \hspace{\stretch{1}}  Pardubice, 2022}
        
            
    }
\end{titlepage}

\cleardoublepage

%% Úvodní stránka s informacemi
{\bfseries
    \begin{center}
        \LARGE{STŘEDOŠKOLSKÁ ODBORNÁ ČINNOST}

        \vspace{14pt}
        {\large
            Obor č. 18: Informatika
        }

        \vspace{0.3 \textheight}

        \LARGE{
        Codeventure
        }

        \LARGE{
        Codeventure
        }

        \vspace{0.24\textheight}
    \end{center}  
}
{\Large
    \noindent\textbf{Autoři:} Michal Maděra a Josef Kahoun\\
    \textbf{Škola:} DELTA - Střední škola informatiky a ekonomie, s.r.o.\\
    \textbf{Kraj:} Pardubický kraj\\
    \textbf{Konzultanti:} RNDr. Jan Koupil, Ph.D. a akademický malíř Daniel Václavík\\
}

\noindent Pardubice, 2022

\cleardoublepage

\noindent{\Large{\bfseries{Prohlášení}}}

\noindent Prohlašujeme, že jsme svou práci SOČ vypracovali samostatně a použili jsme pouze prameny a literaturu uvedené v seznamu bibliografických záznamů.

\noindent Prohlašujeme, že tištěná verze a elektronická verze soutěžní práce SOČ jsou shodné. 

\noindent Nemáme závažný důvod proti zpřístupňování této práce v souladu se zákonem č. 121/2000 Sb., o právu autorském, o právech souvisejících s právem autorským a o změně některých zákonů (autorský zákon) ve znění pozdějších předpisů. 

\vspace{24 pt}

\noindent V Pardubicích dne x. xxxx 2022 \dotfill{}\hspace{\stretch{0.5}} 

\hspace{6cm} Michal Maděra, Josef Kahoun

\cleardoublepage

\vspace*{0.8\textheight}
\noindent{\Large{\bfseries{Poděkování}}}

\noindent
Chtěli bychom poděkovat našim učitelům: RNDr. Janu Koupilovi, Ph.D. za vedení projektu po technické stránce a akademickému malíři Danielu Václavíkovi za vedení práce po grafické stránce. Dále bychom chtěli poděkovat všem lidem, kteří se zapojili testovací fáze našeho projektu.

\cleardoublepage

\noindent{\Large{\bfseries{Abstrakt}}}

\noindent Práce SOČ dokumentuje návrh a vývoj hry Codeventure zaměřené na rozvoj algoritmických a programovacích dovedností u dětí mladšího školního věku. Realizace projektu zahrnuje backendovou část, mobilní aplikaci v roli ovladače hry a webový frontend v roli displeje. Součástí řešení je využití umělé inteligence pro strojové rozpoznání obrazu.

\vspace{18pt}

\noindent{\Large{\bfseries{Klíčová slova}}}

\noindent Programování; GraphQL; TensorFlow 

\vspace{18pt}

\noindent{\Large{\bfseries{Abstract}}}

\noindent The main goal of this project is to create an educational game for primary school students, in which they learn how to logically think about code.

\vspace{18pt}

\noindent{\Large{\bfseries{Keywords}}}

\noindent Programming; GraphQL; TensorFlow

\cleardoublepage

\tableofcontents

\pagenumbering{arabic}
\pagestyle{fancy}
\setcounter{page}{1}

\chapter{Úvod}

Celý svět kolem nás používá počítačové technologie s naprogramovanými aplikacemi. Právě kvůli tomu se s programováním setkáváme každý den, aniž bychom o tom věděli. Jedná se opravdu o užitečnou schopnost, kterou chceme rozvíjet. Samozřejmě je nejlepší začít v raném věku a to u dětí ve školním období. Je rozumné, když se malé děti setkají s programováním co nejdříve. Nejedná se pouze o schopnost umět naprogramovat aplikaci, ale především zapojení a trénování logického myšlení. Nejlepším způsobem jak se těmto schopnostem naučit je jednoznačně formou hry. V dnešní době existují hry, jako LightBot, Kodable a mnoho dalších, které děti rozvíjejí právě ve zmíněných okruzích. Často se jedná o hry, které jsou na menší zařízení, jako tablety, či telefony a dítě prochází hrou pouze za pomocí dotykové obrazovky. Dětem obecně pomáhá, když můžou při učení se nové věci manipulovat s fyzickým předmětem, protože zapojují více smyslů a díky kontaktu a prožitku se učí nové věci více kvalitně a v kratším časovém intervalu. My jsme chtěli do této problematiky také přispět, a proto jsme vymysleli a poté vytvořili hru s názvem Codeventure. Celkový návrh hry je navržený tak, abychom dokázali děti co nejefektivněji naučit zapojit logické myšlení a rozvinuli jejich schopnost základů programování. Hráč pro průchod hrou využívá fyzické kartičky, které představují jednotlivé herní kroky kvůli lepšímu učení. Pro nejlepší herní prožitek je vhodné si hru zapnout na velké zobrazovací ploše a využívat mobilní aplikaci pouze jako ovladač, který slouží pro přepínání a plnění jednotlivých herních úrovní. Hra využívá také umělou inteligenci, která rozpoznává jednotlivé kartičky naskenované hráčem po složení herní sekvence. Díky hře Codeventure dítě zažije nezapomenutelný zážitek, při kterém se naučí využívat logické myšlení, základy programování a ještě se přitom zabavý v kouzelném světě Codeventure, který čeká opravdu na kohokoliv.
\chapter{Teoretická část}

\section{Pojmy}

\subsection{WebSocket}

\subsection{Frontend}

\subsection{Backend}

\subsection{Kompilace}

\subsection{Kontejner}
A container is a standard unit of software that packages up code and all its dependencies so the application runs quickly and reliably from one computing environment to another. \cite{Kontejner}

\subsection{Framework}

\subsection{Komponenta}
Komponenta je menší část aplikace, která poskytuje určitou funkcionalitu aplikace.

\subsection{Herní engine}

\subsection{Cloud}

\subsection{VPS}

\subsection{API}

\subsection{Škálování}

\subsection{JSON}

\subsection{Transpilace}

\subsection{Transakce}

\subsection{Cache}

\section{Technologie}

\subsection{Unity}
Unity je herní engine určený hlavně pro nováčky, ale i pro pokročilé programátory. Jako programovací jazyk Unity využívá C\# vyvýjený firmou Microsoft. Tento programovací jazyk je jednoduchý pro nováčky a je jedním z nejzsnažších programovacích jazyků v současnosti na naučení. \cite{Csharp} Dají se v něm vytvářet 2D i 3D aplikace. Výsledná hra je spustitelná na mnoha platformách, od počítačů až po mobilní zařízení.

\subsection{JavaScript}
JavaScript je moderní programovací jazyk využívající se jak na frontendu, ale i na backendu. Jazyk je tzv. interpretovaný jazyk, což znamená, že kód je spustitelný rovnou ze zdrojového kódu a není potřeba jej kompilovat.

\subsection{TypeScript}
TypeScript je nadstavbou JavaScriptu, která poskytuje typování kódu. Díky tomu je ve výsledné aplikaci méně chyb spojených s typy, protože jsou kontrolovány během vývoje a sestavování aplikace. Ve výsledku je kód transpilován zpět do JavaScriptu.
TODO domyslet

\subsection{React.js}
React.js je JavaScriptový framework pro vývoj uživatelské rozhraní neboli UI. Tvoří se s ním dynamické webové aplikace nebo mobilní aplikace. Vývojářem tohoto frameworku je firma Meta (dříve známá jako firma Facebook). UI je děleno na menší komponenty, které mají svoji funkcionalitu a dají se mezi sebou jednoduše zaměňovat. Díky tomu není například potřeba načíst celou stránku při přechodu na jinou stránku, ale stačí pouze zaměnit komponenty které se na stránce změnily (např hlavička a patička zůstávají stále načteny, ale obsah stránky se změní).

\subsection{Next.js}
Next.js je JavaScriptový framework stavící na frameworku React.js. Je vytvářen společností Vercel, která pod stejným názvem provozuje i vlastní Cloud. Tento framework usnadňuje vývojáři práci se stránkami a přidává nezpočet výhod nad použitím samothného Reactu.js. Hlavním lákadlem je takzvaný "Server Side Rendering" což znamená, že počáteční kód který, je poslán klientovi, je nejdříve vykreslen na serveru, pak je odeslán v neživé podobě klientovi a až po té, co je stránka načtena tak se z něj stává dynamická stránka pomocí React.js. Mezi další výhody patří například automatická optimalizace velikosti obrázků podle toho, jak velké je klient potřebuje, aby se nemusely posílat zbytečně velké a neoptimalizované obrázky.

\subsection{Node.js}
Node.js je backendové prostředí běžící na JavaScriptovém kódu. Výhodou oproti starším programovacím jazykům které běží na serveru např. PHP je, že aplikace neustále běží na pozadí. To umožňuje o dost jednoduší práci s časem (např. načasování určitých událostí neboli Cronjobů) nebo práci v reálném čase pomocí WebSocketů.

\subsection{GraphQL}
GraphQL je je jazyk pro API rozhraní sloužící pro jednoduchou komunikací mezi aplikací a serverem. Typy dotazů se dělí na 3 skupiny: Query: získání dat ze serveru, Mutation: změna dat na serveru a Subscription: zasílání dat v reálném čase pomocí WebSocketů.

\subsection{MongoDB}
MongoDB je dokumentová databáze, která je používaná na velké objemy dat. Data jsou zde ukládána do dokumentů ve formátu JSON. Databáze je jednoduše škálovatelná mezi více zařízeními v tak zvaném "Replica Setu" (data jsou kopírovaná mezi všemi zařízením, aby bylo možné při výpadku některého zařízení pokračovat dál v chodu aplikace). Replica set umožňuje i další pokročilé funkce jako například "transakce".

\subsection{Redis}
Redis je databáze fungující na principu klíč a hodnota. Pod nějakým klíčovým slovem si můžeme uložit hodnotu a následně si můžeme tuto hodnotu zobrazit po zadání tohoto klíče. Databáze funguje primárně v operační paměti počítače, takže funguje velmi rychle a je velmi efektivní. Největší využití je pro ukládání do cache aby byl zrychlen chod aplikace. Databáze také podporuje poslouchání na určitých hodnotách a v reálném čase posílání hodnot na poslouchající klienty. Tato funkcionalita je nejideálnější s technologií WebSocket, nebo ideálněji s GraphQL komponentou "Subscription". 

\subsection{Docker}

\subsection{Kubernetes}
Kubernetes, also known as K8s, is an open-source system for automating deployment, scaling, and management of containerized applications. \cite{Kubernetes}

\subsection{TensorFlow}

\chapter{Tipy k psaní}

Jak už jsem psal výše \LaTeX je dosti komplexní systém, který umožňuje psát velmi rozsáhlé text. Jeho autor Donald Knuth ho stvořil, aby mohl vydat jeho učebnici \emph{The Art of Computer Programming} a dodnes se je využíván pro sazbu skript, učebnic, článků či závěrečných prací. V této kapitole najdeš ukázky různých funkcí a balíčků \LaTeX u od těch nejzákladnějších až po složitější. Neznamená to nutně, že všechny musíš použít, ale když potřebuješ pomoct, tak je dobré mít oporu. 

Pokud s \LaTeX em úplně začínáš tak ti můžu doporučit přiručku \emph{Ne příliš stručný úvod do systému \LaTeX2e}~\cite{LaTeXprirucka}. Případně spoustu užitečných informací nalezneš na Wikibooks~\cite{wikibooksLaTeX}. Pokud narazíš na nějaký problém googli. Na internetu je spoustu fór, kde pravděpodobně už někdo podobný problém řešil. Asi nejvíce otho najdeš na stránce \emph{TeX - LaTeX Stackexchange} \cite{stackExchange}.


\section[Základy]{Základy: Text, obrázky, tabulky a citace} %%[Text, který bude v obsahu]{Text, který se vytiskne na stránce} Zkus měnit jednotlivé závorky a uvidíš :) 
Psaní v \LaTeX{u} není žádná věda, stačí psát normálně do zdrojového souboru. Pokud bys chtěl psát obrážky či číslovaný seznam, pak můžeš použít prostředí \texttt{itemize} či \texttt{enumerate}. Často je důležité používat nezlomitelnou mezeru. Tu uděláš pomocí \verb|~|~(tildy). Pokud budeš chtít psát uvozovky použij příkaz \texttt{uv}, pomocí něj se ti vytvoří uvozovky podle příslušného jazyka. V česku tedy ve formátu 99 66. Použití příkazu najdeš níže v textu.

Občas je zapotřebí \LaTeX{u} pomoct při rozdělování slov. To se udělá snadno vložením symbolů \verb|\-| mezi jednotlivé slabiky.

\subsection{Obrázky}

U obrázků je dobré používat vektorové formáty, pokud to jde. \LaTeX se nejvíc kamarádí s formátem PDF. Do známého PDFka lze z jiných vektorových formátů (ať už SVG či ESP) obrázky přenést snadno pomocí grafických programů, jako je třeba Inkscape. \LaTeX si rozhodně poradí i s tradičními formáty PNG a JPG, avšak tyto obrázky mohou zabírat více prostoru a při tisku se může projevit nižší rozlišení obrázků. Pokud chceš používat tyto obrázky, rozhodně měj na paměti, aby měli rozlišení alespoň 250 indálně 330 ppi.

Obrázky se vkládají do prostředí \texttt{figure}, při úpravě šířky je možné krom tradičních jednotek jako cm nebo mm použít také jako jednotku šířku stránky \texttt{textwidth} to se hodí zejména když chceš mít více podobrázků. 

U každého obrázku je důležité aby měl popisek, \texttt{caption}. Do popisku napiš, co na obrázku je, případně nějaký další popis, tak aby čtenář následně neměl sebemenší pochybnost. U obrázků co nejsou tvoje nezapomeň an citaci. Jinak by to totiž znamenalo, že jsi obrázek dělal ty sám, což není etické přivlastňovat si cizí díla. Popisek obrázku je věta, proto musí vždy končit tečkou.

\begin{figure}
    \centering %% příkaz, který ti obrázek zarovná na střed
    \includegraphics[width=0.6\textwidth]{img/soc-logo.jpg} %% vložení samotného obrátku
    \caption{Logo SOČky \cite{socLogo}.} %% popisek obrázku, nezapomeň na citace!
    \label{fig:logoSOC} %% označení až budeš chtít na obrázek odkazovat
\end{figure}

Když chceš odkazovat na obrázek, stačí pak už jen napsat příkaz \texttt{ref} a do závorek napsat označení obrázku. Třeba logo SOČky, můžeš vidět na obrázku \ref{fig:logoSOC} \cite{socLogo}.

Pokud bys měl více podobrázků přichází do hry balíček \texttt{subcaption}. Pomocí něj lze vysázet i podobrázky. U podobrázků se popisek píše pouze jeden, dolů. Je v tomto připadě vhodné použít navíc hranaté závorky, do nichž se napíše kratší popisek, který se následně ukáže v seznamu obrázků.


\begin{figure}[h] \centering
    % \begin{subfigure}{0.63\textwidth} %%prostředí pro podobrázek {šířka podobrázku}
    %   \includegraphics[width=\textwidth]{imgs/vysledek_10} 
    %   \caption{} %% aby se ti vysázelo označení obrázku.
    % \end{subfigure}
    % \begin{subfigure}{0.63\textwidth}
    %   \includegraphics[width=\textwidth]{imgs/vysledek_20}
    %   \caption{}
    % \end{subfigure}
    \caption[Graf závislosti rotace DH PSF $\Delta\varphi$ na defokusaci objektivu $\Delta z$.]{Graf závislosti rotace DH PSF $\Delta\varphi$ na defokusaci objektivu $\Delta z$, (a) při použití objektivu Plan Fluor 10, (b) při použití objektivu Plan Fluor 20. Měřená data (žluté body) jsou lineárně proloženy (přerušovaná přímka). }
    \label{fig:rotace_grafy}
\end{figure}

Všimni si, že obrázky jsou naschvál široké. Je to proto, aby byly dobře čitelné. Také si všimni popisku grafů. Ačkoli nejspíš netušíš co je to DH PSF či defokusace objektivu mělo by ti být jasné, že je důležité přesně graf popsat. To znamená co je na vodorovné ose, co je na svislé ose. V jakých jednotkách veličiny jsou. Které body co znamenají, která křivka má jaký význam. Napsat samotné \uv{$\Delta \varphi$} je málo, vždy raději připoměň, co daná značka znamená.

\subsection{Tabulky}

U tabulek platí to stejné co u obrázků. Zarovnávají se na střed a nechávají se \uv{plavat} v textu. Tabulka narozdíl od textu, má popisek nahoře. U tabulky \ref{tab:ukazka} je použit balíček \texttt{booktabs}, pomocí kterého je celá tabulka naformátovaná.

Seznam jak obrázků tak tabulek je pak vytvořen pomocí příkazů \texttt{listoftables} a~\texttt{list\-of\-fig\-ures} na konci práce před literaturou.

\begin{table}[b]
    \caption{Tato tabulka slouží jako ukázka toho, jak mohou tabulky vypadat.} %% popisek se u tebaluky píše nad ní
    \label{tab:ukazka} %% označení pro pozdější odkazování se 
    \centering
        \begin{tabular}{lll}
            \toprule %% příkazy z balíčku booktabs
            záhlaví& této & tabulky\\
            \midrule
            obsah&tabulky& už\\
            není & oddělený &čarami\\
            \bottomrule
        \end{tabular}
\end{table}

\subsection{Literatura}

V \LaTeX{}u lze dělat seznam literatury dvěma způsoby. V této šabloně jsem použil ten, kdy se seznam literatury píše přímo do práce. Pro jeho vygenerování doporučuji použít některý z generátarů, jako jsou například Citace PRO \cite{citacePRO}. Pomocí citací lze vygenerovat přímo dokument, který se pak už jen překopíruje do textu a člověk nemusí nic zvýrazňovat. Dále lze využít Bibtex, který rozhodně do budoucna hodlám zaimplementovat do šablony, avšak jeho použití nemusí být tak přátelské k začátečníkům.

Pokud bys chtěl odkazovat na vícero zdrojů stačí je napsat vedle sebe oddělené čárkou \cite{LaTeXprirucka, citacePRO, Born2019}. Případně můžu odkaz na konkrétní stránku dát do hranatých závorek, viz \cite[str.~1]{Born2019}


\section[Pokročilejší tipy]{Pokročilejší tipy, které se mohou hodit}

\subsection{Rovnice}

Jak můžeš vidět tak rovnice lze psát jednak do textu a nebo pokud se jedná o nějakou důležitou nebo rozsáhlejší rovnici tak na samostatný řádek. Pokud je rovnice opravdu důležitá, tak je vhodné ji také číslovat. Pak se na ni můžeš dále odkazovat v textu.
\begin{equation}
    \vec{F} = m \vec{a}
    \label{eq:newton2}
\end{equation}
\dots Například podle druhého Newtonova zákona, rovnice (\ref{eq:newton2}) \dots Zároveň je vždy nutné vysvětlit co která veličina znamená. V tomto případě bych napsal, že v druhém Newtonově zákoně vektor síly $\vec F$ odpovídá součinu hmotnosti tělesa $m$ a jeho zrychlení $\vec a$. 

Věřím, že se sazbou matematiky ti pomůže tvůj školitel, případně mi můžeš napsat (mail je v úvodu). Jednotlivé funkcionality spolu se seznamem znaků nalezneš jednak v Ne příliš stručném úvodu~\cite{LaTeXprirucka} nebo na Wikibooks v sekcích \emph{Mathematics} a \emph{Advanced mathematics}~\cite{wikibooksLaTeX}.

\chapter{Když dokončuji práci}

Každou práci je dobré zkontrolovat, aby v ní nebyly pravopisné chyby, nebyla těžkopádně napsaná -- byla čtivá -- a neobsahovala žádný typografický nedostatek. Proto, když práci sepíšeš, nech ji chvilku odležet, třeba týden. Pak si ji po sobě znovu přečti. Hned uvidíš, kolik věcí bys napsal jinak případně kde tě bije do očí jaká chyba. Dej práci přečíst také svému školiteli a případně češtináři. Zajistíš tak, že bude obsahovat méně chyb.

Pak můžeš práci vytisknout a hurá do soutěže.

\chapter{Závěr}

Věřím, že jsem ti spolu se šablonou poskytl několik tipů, jak napsat práci. Ať už jde o úplné začátky s \LaTeX{}em. Či ukázku toho, co vše s ním zvládneš. Pokud bys měl k šabloně libovolné dotazy, rouhodně se na mě obrať. \LaTeX tvé práci dodá určitou krásu, tak doufám, že ti dodá sebevědomí a uspěješ při souteži. A i kdyby ne vzpomeň si, kolik ses toho musel naučit a hned uvidíš o jaký kus ses posunul.

%% literatura
\begin{thebibliography}{99}
    \bibitem{Kontejner} \textit{Kontejner} [online]. [cit. 2022-01-10]. Dostupné z \url{https://www.docker.com/resources/what-container/}
    \bibitem{Framework} \textit{Framework} [online]. [cit. 2022-01-22]. Dostupné z \url{https://codeinstitute.net/global/blog/what-is-a-framework/}
    \bibitem{Skalovani} \textit{Škálování} [online]. [cit. 2022-01-22]. Dostupné z \url{https://www.cyberlinkasp.com/insights/what-is-software-scalability-and-why-is-it-important/}
    \bibitem{JSON} \textit{JSON} [online]. [cit. 2022-01-22]. Dostupné z \url{https://www.json.org/json-en.html}
    \bibitem{Transpilace} \textit{Transpilace} [online]. [cit. 2022-01-22]. Dostupné z \url{https://www.digitalocean.com/community/tutorials/javascript-transpilers-what-they-are-why-we-need-them}
    \bibitem{Cache} \textit{Cache} [online]. [cit. 2022-01-22]. Dostupné z \url{https://www.businessinsider.com/what-is-cache}
    \bibitem{Flutter} \textit{Flutter} [online]. [cit. 2022-01-22]. Dostupné z \url{https://flutter.dev/}
    \bibitem{OpenSource} \textit{Open Source} [online]. [cit. 2022-01-22]. Dostupné z \url{https://opensource.com/resources/what-open-source}
    \bibitem{Csharp} \textit{Nejjednoduší programovací jazyky na naučení} [online]. [cit. 2021-01-10]. Dostupné z \url{https://careerkarma.com/blog/easiest-programming-languages-to-learn/#Csharp}
    \bibitem{TypeScript} \textit{TypeScript} [online]. [cit. 2022-01-10]. Dostupné z \url{https://www.typescriptlang.org/}
    \bibitem{Prisma} \textit{Prisma} [online]. [cit. 2022-01-28]. Dostupné z \url{https://www.prisma.io/docs/concepts/overview/what-is-prisma}
    \bibitem{PubSub} \textit{PubSub} [online]. [cit. 2022-01-28]. Dostupné z \url{https://redis.io/topics/pubsub}
    \bibitem{bcrypt} \textit{bcrypt} [online]. [cit. 2022-01-28]. Dostupné z \url{https://auth0.com/blog/hashing-in-action-understanding-bcrypt/}
    \bibitem{CDN} \textit{CDN} [online]. [cit. 2022-02-12]. Dostupné z \url{https://www.cloudflare.com/learning/cdn/what-is-a-cdn/}
    \bibitem{Cloudflare-pages} \textit{Cloudflare Pages} [online]. [cit. 2022-02-12]. Dostupné z \url{https://pages.cloudflare.com/}
    \bibitem{Cloudflare-cdn} \textit{Cloudflare CDN} [online]. [cit. 2022-02-12]. Dostupné z \url{https://www.cloudflare.com/cdn/}
    \bibitem{react-unity-webgl} \textit{react-unity-webgl} [online]. [cit. 2022-02-12]. Dostupné z \url{https://www.npmjs.com/package/react-unity-webgl}
    \bibitem{AssetBundle} \textit{AssetBundle} [online]. [cit. 2022-02-13]. Dostupné z \url{https://docs.unity3d.com/Manual/AssetBundlesIntro.html}
    \bibitem{Coroutine} \textit{Coroutine} [online]. [cit. 2022-02-13]. Dostupné z \url{https://docs.unity3d.com/Manual/Coroutines.html}
    \bibitem{SharedPreferences} \textit{SharedPreferences} [online]. [cit. 2022-02-13]. Dostupné z \url{https://developer.android.com/training/data-storage/shared-preferences}
    \bibitem{NSUserDefaults} \textit{NSUserDefaults} [online]. [cit. 2022-02-13]. Dostupné z \url{https://developer.apple.com/documentation/foundation/nsuserdefaults}
    \bibitem{ImagePicker} \textit{ImagePicker} [online]. [cit. 2022-02-13]. Dostupné z \url{https://pub.dev/packages/image_picker}
    \bibitem{OWASP} \textit{OWASP} [online]. [cit. 2022-03-03]. Dostupné z \url{https://owasp.org/}
    \bibitem{OWASP-top-ten} \textit{OWASP Top 10} [online]. [cit. 2022-03-03]. Dostupné z \url{https://owasp.org/www-project-top-ten/}
    \bibitem{CryptographicFailures} \textit{Kryptografické selhání} [online]. [cit. 2022-03-03]. Dostupné z \url{https://owasp.org/Top10/A02_2021-Cryptographic_Failures/}
    \bibitem{Injection} \textit{Injection} [online]. [cit. 2022-03-03]. Dostupné z \url{https://owasp.org/Top10/A03_2021-Injection/}
    \bibitem{Validator} \textit{Validator} [online]. [cit. 2022-03-03]. Dostupné z \url{https://www.npmjs.com/package/validator}
    \bibitem{VulnerableAndOutdatedComponents} \textit{Vulnerable and Outdated Components} [online]. [cit. 2022-03-03]. Dostupné z \url{https://owasp.org/Top10/A06_2021-Vulnerable_and_Outdated_Components/}
    \bibitem{What-is-AI} \textit{What is AI} [online]. [cit. 2022-03-04]. Dostupné z \url{https://borghese.di.unimi.it/Teaching/AdvancedIntelligentSystems/Old/IntelligentSystems_2008_2009/Old/IntelligentSystems_2005_2006/Documents/Symbolic/04_McCarthy_whatisai.pdf}
    \bibitem{IBM-AI} \textit{IBM - What is AI} [online]. [cit. 2022-03-04]. Dostupné z \url{https://www.ibm.com/cloud/learn/what-is-artificial-intelligence}
    \bibitem{IBM-CV} \textit{IBM - What is computer vision} [online]. [cit. 2022-03-04]. Dostupné z \url{https://www.ibm.com/topics/computer-vision}
\end{thebibliography}


%% obrázky 
\listoffigures

%% tabulky
\listoftables

\appendix %% začínají přílohy

\titleformat{\chapter}[block]{\scshape\bfseries\LARGE}{Příloha \thechapter}{10pt}{\vspace{0pt}}[\vspace{-22pt}] %% nastavení nadpisu u příloh


\chapter{%Příloha A 
Spot diagramy a další }


\end{document}